\documentclass[11pt,a4paper]{article}

\usepackage[a4paper, margin=2.5cm]{geometry}
\usepackage{booktabs}
\usepackage{array}
\usepackage{longtable}
\usepackage{hyperref}
\usepackage{xcolor}
\usepackage{titlesec}
\usepackage{parskip}
\usepackage{fancyhdr}
\usepackage{courier}
\usepackage{listings}

\definecolor{codegray}{rgb}{0.95,0.95,0.95}
\lstset{
  backgroundcolor=\color{codegray},
  basicstyle=\ttfamily\small,
  frame=single,
  framesep=4pt,
  breaklines=true
}

\hypersetup{colorlinks=true, linkcolor=blue, urlcolor=blue}

\pagestyle{fancy}
\fancyhf{}
\rhead{CIF Microarchitecture}
\lhead{Transaction Tag Allocator}
\rfoot{\thepage}

\title{\textbf{CIF Microarchitecture Specification}\\[0.4em]
\large Block: Transaction Tag Allocator}
\author{CIF Design Team}
\date{\today}

\begin{document}

\maketitle
\thispagestyle{fancy}
\tableofcontents
\newpage

% -----------------------------------------------------------------------
\section{Overview}

The Transaction Tag Allocator is a shared combinational/registered block
instantiated once in the CIF and used by both the read and write paths
across all \texttt{N\_CLIENTS = 32} AXI clients.

Its responsibilities are:
\begin{itemize}
  \item Allocate a unique \texttt{txn\_tag[7:0]} for every validated
        transaction entering the pipeline.
  \item Enforce the per-AXID outstanding transaction limit of
        \texttt{MAX\_OUTSTANDING = 4}.
  \item Stall upstream (Request Validator / Error Handler) on a per-AXID
        basis when the outstanding limit is reached.
  \item Free an occupied slot when a transaction completes — via
        \texttt{BRESP} on the write path, or \texttt{rob\_last} on the
        read path.
\end{itemize}

Tag encoding (agreed with MC Core team):
\begin{lstlisting}
txn_tag[7:0]:
  [7:2] = AXID[5:0]   -- 6-bit AXI ID, up to 64 unique IDs
  [1:0] = seqnum[1:0] -- per-AXID slot index, wraps at 4
\end{lstlisting}

% -----------------------------------------------------------------------
\section{Design Parameters}

\begin{table}[h!]
\centering
\begin{tabular}{lll}
\toprule
\textbf{Parameter} & \textbf{Value} & \textbf{Description} \\
\midrule
\texttt{AXID\_W}           & 6   & AXI ID width; 64 unique IDs \\
\texttt{MAX\_OUTSTANDING}  & 4   & Max outstanding txns per AXID \\
\texttt{SEQNUM\_W}         & 2   & Seqnum width (\(\log_2\) MAX\_OUTSTANDING) \\
\texttt{TXN\_TAG\_W}       & 8   & Tag width = AXID\_W + SEQNUM\_W \\
\texttt{N\_AXIDS}          & 64  & Total AXID space ($2^6$) \\
\texttt{OCC\_CTR\_W}       & 2   & Occupancy counter width per AXID \\
\bottomrule
\end{tabular}
\caption{Transaction Tag Allocator parameters}
\end{table}

% -----------------------------------------------------------------------
\section{Internal Blocks}

\subsection{Occupancy Counter Array}

\begin{itemize}
  \item \textbf{Structure:} 64 independent 2-bit saturating counters,
        one per AXID.
  \item \textbf{Increment:} On every successful allocation
        (\texttt{alloc\_en} asserted), the counter for the requesting
        AXID is incremented.
  \item \textbf{Decrement:} On assertion of \texttt{WR\_FREE\_VALID} or
        \texttt{RD\_FREE\_VALID}, the counter for the AXID extracted from
        the incoming tag (\texttt{tag[7:2]}) is decremented.
  \item \textbf{Read port:} Continuously exposes
        \texttt{occupancy[AXID]} to Stall Logic.
\end{itemize}

\subsection{Seqnum Counter Array}

\begin{itemize}
  \item \textbf{Structure:} 64 independent 2-bit wrap counters, one per
        AXID. Represents the \emph{next} seqnum to be allocated.
  \item \textbf{Increment:} On every successful allocation
        (\texttt{alloc\_en} asserted), the counter for the requesting
        AXID is incremented modulo 4.
  \item \textbf{Read port:} Continuously exposes \texttt{seqnum[1:0]}
        for the requesting AXID to Allocate Logic.
  \item \textbf{No decrement:} Seqnum counters are allocation-only;
        they are never decremented. The occupancy counter separately
        tracks the number of in-flight slots.
\end{itemize}

\subsection{Stall Logic}

\begin{itemize}
  \item \textbf{Type:} Purely combinational.
  \item \textbf{Function:} Asserts \texttt{STALL} for a given AXID when
        \texttt{occupancy[AXID] == MAX\_OUTSTANDING (4)}.
  \item \textbf{Granularity:} Per-AXID. Only the AXID currently
        presented by the upstream block is evaluated; other AXIDs are
        unaffected.
  \item \textbf{Output:} \texttt{STALL} fed back to Request Validator
        and Error Handler.
\end{itemize}

\subsection{Allocate Logic}

\begin{itemize}
  \item \textbf{Gate condition:} Fires only when \texttt{REQ\_VALID}
        is asserted and \texttt{STALL} is deasserted
        (\texttt{alloc\_en = REQ\_VALID \& \textasciitilde{}STALL}).
  \item \textbf{Tag formation:}
\begin{lstlisting}
TXN_TAG[7:0] = { AXID[5:0], seqnum[1:0] }
\end{lstlisting}
  \item \textbf{Side-effects on alloc:}
        \texttt{inc\_occ} to Occupancy Counter Array;
        \texttt{inc\_seq} to Seqnum Counter Array.
  \item \textbf{Output:} \texttt{TXN\_TAG[7:0]} and \texttt{TAG\_VALID}
        forwarded to Burst Splitter.
  \item \textbf{Error path:} Error Handler presents AXID with
        \texttt{REQ\_VALID} asserted; allocation proceeds identically.
        The resulting tag is used to form a NOP packet
        (\texttt{pkt\_type = Err\_write / Err\_read}) in the normal
        pipeline. No special casing in Allocate Logic.
\end{itemize}

\subsection{Free Logic}

\begin{itemize}
  \item \textbf{Two free ports:}
  \begin{itemize}
    \item Write path: \texttt{WR\_FREE\_VALID}, \texttt{WR\_FREE\_TAG[7:0]}
          — sourced from Response Tracker on BRESP emission.
    \item Read path: \texttt{RD\_FREE\_VALID}, \texttt{RD\_FREE\_TAG[7:0]}
          — sourced from Read Transaction Table on \texttt{rob\_last}
          assertion.
  \end{itemize}
  \item \textbf{AXID extraction:} \texttt{AXID = tag[7:2]} in both cases.
  \item \textbf{Action:} Asserts \texttt{dec\_occ} to Occupancy Counter
        Array for the extracted AXID.
  \item \textbf{Simultaneous frees:} Both ports may assert in the same
        cycle for different AXIDs — handled independently. Simultaneous
        frees on the same AXID (one write + one read) are not possible
        by construction (a tag is either a read or a write transaction,
        never both).
\end{itemize}

% -----------------------------------------------------------------------
\section{Interface Signals}

\subsection{Inputs}

\begin{longtable}{llll}
\toprule
\textbf{Signal} & \textbf{Width} & \textbf{Source} & \textbf{Description} \\
\midrule
\texttt{REQ\_VALID}       & 1 & Req Validator / Error Handler & Allocation request valid \\
\texttt{AXID[5:0]}        & 6 & Req Validator / Error Handler & Requesting AXI ID \\
\texttt{WR\_FREE\_VALID}  & 1 & Response Tracker              & Write path free strobe \\
\texttt{WR\_FREE\_TAG[7:0]} & 8 & Response Tracker            & Tag being freed (write) \\
\texttt{RD\_FREE\_VALID}  & 1 & Read Txn Table                & Read path free strobe \\
\texttt{RD\_FREE\_TAG[7:0]} & 8 & Read Txn Table              & Tag being freed (read) \\
\bottomrule
\caption{Tag Allocator input signals}
\end{longtable}

\subsection{Outputs}

\begin{longtable}{llll}
\toprule
\textbf{Signal} & \textbf{Width} & \textbf{Destination} & \textbf{Description} \\
\midrule
\texttt{TXN\_TAG[7:0]} & 8 & Burst Splitter                & Allocated transaction tag \\
\texttt{TAG\_VALID}    & 1 & Burst Splitter                & Tag output valid \\
\texttt{STALL}         & 1 & Req Validator / Error Handler & Per-AXID stall signal \\
\bottomrule
\caption{Tag Allocator output signals}
\end{longtable}

% -----------------------------------------------------------------------
\section{Functional Behaviour}

\subsection{Allocation Flow}

\begin{enumerate}
  \item Upstream presents \texttt{AXID[5:0]} and asserts
        \texttt{REQ\_VALID}.
  \item Stall Logic reads \texttt{occupancy[AXID]} from the Occupancy
        Counter Array combinationally.
  \item If \texttt{occupancy[AXID] < 4}: \texttt{STALL} is deasserted,
        \texttt{alloc\_en} is asserted.
  \item Allocate Logic reads \texttt{seqnum[1:0]} from the Seqnum
        Counter Array for the requesting AXID.
  \item \texttt{TXN\_TAG = \{AXID, seqnum\}} is driven; \texttt{TAG\_VALID}
        is asserted.
  \item Occupancy Counter Array increments for AXID.
  \item Seqnum Counter Array increments (mod 4) for AXID.
\end{enumerate}

\subsection{Stall Flow}

\begin{enumerate}
  \item If \texttt{occupancy[AXID] == 4}: \texttt{STALL} is asserted
        to upstream.
  \item Upstream holds \texttt{REQ\_VALID} and \texttt{AXID} stable.
  \item Stall is released as soon as Free Logic decrements occupancy
        for that AXID (i.e., on the next \texttt{WR\_FREE\_VALID} or
        \texttt{RD\_FREE\_VALID} for a tag belonging to that AXID).
\end{enumerate}

\subsection{Free Flow}

\begin{enumerate}
  \item Response Tracker asserts \texttt{WR\_FREE\_VALID} with
        \texttt{WR\_FREE\_TAG} when a BRESP is emitted (write path).
  \item Read Txn Table asserts \texttt{RD\_FREE\_VALID} with
        \texttt{RD\_FREE\_TAG} when \texttt{rob\_last} fires (read path).
  \item Free Logic extracts \texttt{AXID = tag[7:2]} and asserts
        \texttt{dec\_occ} to the Occupancy Counter Array.
  \item This may immediately release a stall on that AXID if
        Stall Logic observes the updated occupancy.
\end{enumerate}

% -----------------------------------------------------------------------
\section{Timing Notes}

\begin{itemize}
  \item Stall Logic is purely combinational: \texttt{STALL} is available
        in the same cycle that \texttt{AXID} is presented.
  \item Allocate Logic is registered: \texttt{TXN\_TAG} and
        \texttt{TAG\_VALID} are valid on the cycle \emph{following}
        \texttt{alloc\_en}.
  \item Counter updates (inc/dec) are registered and take effect on the
        following clock edge. A free in cycle $N$ releases the stall in
        cycle $N{+}1$.
\end{itemize}

% -----------------------------------------------------------------------
\section{Edge Cases}

\begin{itemize}
  \item \textbf{Simultaneous alloc and free on same AXID:} If
        \texttt{alloc\_en} and \texttt{dec\_occ} fire in the same cycle
        for the same AXID, occupancy is held constant (net zero change).
        This is a valid steady-state operating condition.
  \item \textbf{Error path:} Error Handler allocates a tag identically
        to valid transactions. The resulting tag is embedded in a NOP
        packet (\texttt{pkt\_type = Err\_write / Err\_read}). The slot
        is freed normally on MC Core error completion.
  \item \textbf{Underflow protection:} The Occupancy Counter Array must
        not decrement below zero. This is guaranteed by construction —
        a free is only issued when a tag is in-flight — but a
        design-time assertion should be included in the verification
        environment.
\end{itemize}

\end{document}

